\documentclass[11pt]{article}
\usepackage{amsmath,amssymb,amsthm,mathtools}
\usepackage[margin=1in]{geometry}
\usepackage{setspace}
\usepackage{hyperref}
\usepackage{algorithm}
\usepackage{algpseudocode}
\setstretch{1.15}
\newtheorem{proposition}{Proposition}[section]
\newtheorem{theorem}[proposition]{Theorem}
\newcommand{\E}{\mathbb E}
\newcommand{\Pp}{\mathbb P}
\newcommand{\R}{\mathbb R}
\newcommand{\Law}{\mathcal L}
\newcommand{\Var}{\operatorname{Var}}
\newcommand{\Cov}{\operatorname{Cov}}
\newcommand{\ind}{\mathbf 1}
\title{A Noise-Transport Update for the SV Model\\
\large Working Draft}
\author{Wells Guan}
\date{November 23, 2026}

\begin{document}
\maketitle

\noindent
The idea is to use the known observation-noise distribution to construct
a nonlinear update of a forecast pair $(H^f,Y^f)$. The construction below
gives an exact posterior draw if the conditional noise CDF is known.
For an empirical predictive distribution, it reduces to inverse-CDF
resampling. For a continuous predictive distribution, computing the
conditional noise CDF remains a numerical step.

\section{The one-step filtering problem}

Consider the SV model
\begin{align*}
H_t&=\mu_h+\phi(H_{t-1}-\mu_h)+qZ_{1,t},\\
R_t&=\mu_r+\exp(H_t/2)Z_{2,t},
\end{align*}
where $|\phi|<1$, $q>0$, and the noises are mutually independent
standard normal variables, independent of the initial state.
The parameters are treated as known.

Fix the past observations $r_{1:t-1}$ and denote the predictive law by
\[
\mu(dh)=\Pp(H_t\in dh\mid R_{1:t-1}=r_{1:t-1}).
\]
In the following, the past is fixed and the time index is omitted.
We allow $\mu$ to be any probability measure on $\R$.

For an observed return $r\ne\mu_r$, set $y=\log(r-\mu_r)^2$. The
transformed observation model is
\[
Y=H+\eta,\qquad \eta=\log Z^2,\qquad Z\sim N(0,1),\quad \eta\perp H.
\]
Its noise CDF and density are
\begin{equation}
F_\eta(u)=2\Phi(e^{u/2})-1,\qquad
g(u)=\frac{1}{\sqrt{2\pi}}
\exp\left(\frac u2-\frac{e^u}{2}\right).
\label{eq:noise}
\end{equation}
The CDF follows from $\Pp(\log Z^2\le u)=\Pp(|Z|\le e^{u/2})$;
differentiating gives the density.

In this model, the conditional density of $R$ given $H=h$ depends on
$r$ only through $(r-\mu_r)^2$. Thus the transformed observation retains
the information about $H$ in a nonzero centered return. Indeed,
\[
g\bigl(\log(r-\mu_r)^2-h\bigr)
=|r-\mu_r|\,p_R(r\mid h),
\]
and the factor on the right does not depend on $h$.
An exactly zero centered return is a null event in the continuous model.
Rounded or censored observations would need their own observation model.

We would like to construct a function $f_y$ such that
\begin{equation}
(H^f,Y^f)\sim\Law(H,Y)
\quad\Longrightarrow\quad
\Law\bigl(f_y(H^f,Y^f)\bigr)=\Law(H\mid Y=y).
\label{eq:target}
\end{equation}
Here $(H^f,Y^f)$ is an auxiliary simulated pair. Its noise
$\eta^f=Y^f-H^f$ is independent of the actual observation, conditional
on the fixed past. This is different from the noise $\eta$ that
generated the actual $Y$: $\eta$ and $Y=H+\eta$ are not independent.

\section{Why an affine update is only a starting point}

Assume finite second moments and denote
\[
m=\E H,\quad P=\Var(H),\qquad
b=\E\eta,\quad V=\Var(\eta)=\frac{\pi^2}{2}.
\]
The usual population EnKF-type update is
\[
H^a_{\mathrm{lin}}=H^f+K(y-Y^f),\qquad K=\frac{P}{P+V}.
\]
Its mean is the best affine estimate
\[
L(y)=m+K(y-m-b).
\]
If
\[
U=H^f-m-K(Y^f-m-b),
\]
then $H^a_{\mathrm{lin}}=L(y)+U$. But the actual conditional law satisfies
\[
\Law(H\mid Y=y)=L(y)+\Law(U\mid Y^f=y),
\]
where addition of a constant denotes a translation of a distribution.
The affine update instead uses $L(y)+\Law(U)$. Orthogonality of
$U$ and $Y^f$ alone does not identify these two laws.

For a jointly Gaussian pair, the residual is independent of the
observation and this update is exact. For the transformed SV model,
the log-$\chi_1^2$ observation noise is non-Gaussian, so this independence
does not generally hold. We therefore look for a nonlinear update
using the observation-noise distribution.

\section{The conditional noise law}

The identity $H=Y-\eta$ suggests first finding $\Law(\eta\mid Y=y)$.
Define the normalizing factor
\[
c_\mu(y)=\int_\R g(y-h)\,\mu(dh).
\]
The density $g$ is strictly positive and bounded:
$0<g(u)\le (2\pi e)^{-1/2}$. Thus
$0<c_\mu(y)<\infty$ for every $y\in\R$ and every probability measure $\mu$.

\begin{proposition}
A version of the posterior law of $H$ is
\[
\pi_y(A)=\frac{\int_A g(y-h)\,\mu(dh)}{c_\mu(y)}.
\]
The corresponding conditional noise law is
\begin{equation}
\nu_y(B)
=\frac{\int_\R
\ind_{\{y-h\in B\}}g(y-h)\,\mu(dh)}{c_\mu(y)}.
\label{eq:conditional-noise}
\end{equation}
In particular, $\nu_y$ is the image of $\pi_y$ under $h\mapsto y-h$.
\end{proposition}

\begin{proof}
The joint law of $(H,Y)$ is $\mu(dh)g(y-h)\,dy$, which gives the
formula for $\pi_y$. To check the formula for $\nu_y$ directly, let
$B,C$ be Borel subsets of $\R$. By Fubini's theorem,
\begin{align*}
\int_C \nu_y(B)c_\mu(y)\,dy
&=\int_\R\int_C
\ind_{\{y-h\in B\}}g(y-h)\,dy\,\mu(dh)\\
&=\int_\R\int_\R
\ind_{\{u\in B,\ h+u\in C\}}g(u)\,du\,\mu(dh)\\
&=\Pp(\eta\in B,Y\in C).
\end{align*}
This verifies the conditional distribution identity. The relation
$H=y-\eta$ after conditioning gives the final statement.
\end{proof}

Denote the conditional noise CDF by
\begin{equation}
G_y(u)=\nu_y((-\infty,u])
=\frac{\int_\R
\ind_{\{y-h\le u\}}g(y-h)\,\mu(dh)}{c_\mu(y)}.
\label{eq:conditional-cdf}
\end{equation}
This formula is valid even when the predictive law has atoms.

If $\mu$ has a density $p_f$, we may instead write
\[
G_y(u)=
\frac{\int_{-\infty}^u g(v)p_f(y-v)\,dv}
{\int_\R g(v)p_f(y-v)\,dv}.
\]
For example, if $H^f\sim N(m,P)$ with $P>0$, then
\[
p_{\eta\mid Y=y}(u)
\propto
\exp\left\{-\frac{(y-u-m)^2}{2P}
+\frac u2-\frac{e^u}{2}\right\}.
\]
Its log density has second derivative $-1/P-e^u/2<0$.
Thus it is strictly log-concave, although it is generally not Gaussian.
Knowing $g$ does not make $G_y$ equal to $F_\eta$; the predictive law
still enters through $p_f(y-u)$.

\section{An exact nonlinear update}

For $v\in(0,1)$, let
\[
G_y^{-1}(v)=\inf\{u\in\R:G_y(u)\ge v\}.
\]
This generalized inverse also applies when $G_y$ has jumps. Define
\begin{equation}
T_y(u)=G_y^{-1}(F_\eta(u)).
\label{eq:noise-map}
\end{equation}

\begin{theorem}[Noise-transport representation]
Let $H^f\sim\mu$ and let $\eta^f$ be independent of $H^f$ with density $g$.
Set $Y^f=H^f+\eta^f$. For every fixed $y\in\R$, define
\begin{equation}
H^a=f_y(H^f,Y^f),\qquad
f_y(h,z)=y-T_y(z-h).
\label{eq:update}
\end{equation}
Then $\Law(H^a)=\pi_y$, the posterior version specified above.
No Gaussian assumption on $\mu$, and no moment assumption on $\mu$,
is needed for this distributional statement.
\end{theorem}

\begin{proof}
By construction, $Y^f-H^f=\eta^f$. Since $F_\eta$ is continuous,
\[
V^f=F_\eta(Y^f-H^f)=F_\eta(\eta^f)\sim U(0,1).
\]
The inverse-CDF construction gives
\[
T_y(\eta^f)=G_y^{-1}(V^f)\sim\nu_y.
\]
For any Borel set $A$, formula \eqref{eq:conditional-noise} now gives
\begin{align*}
\Pp(H^a\in A)
&=\int_\R\ind_{\{y-u\in A\}}\,\nu_y(du)\\
&=\frac{\int_\R
\ind_{\{h\in A\}}g(y-h)\,\mu(dh)}{c_\mu(y)}
=\pi_y(A),
\end{align*}
and we are done.
\end{proof}

For the SV observation noise, the map is explicitly
\[
f_y(h,z)
=y-G_y^{-1}\left(2\Phi(e^{(z-h)/2})-1\right).
\]
It can also be expressed as a correction of the forecast state:
\[
H^a
=H^f+(y-Y^f)
+\left\{\eta^f-T_y(\eta^f)\right\}.
\]
The last term adjusts the original noise into the noise distribution
compatible with the actual observation.

Notice that this particular construction depends on the forecast pair
only through its noise $\eta^f$. Conditional on the supplied $\mu$ and
the fixed $y$, the resulting $H^a$ is independent of $H^f$, since $\eta^f$
is independent of $H^f$. Thus this is a way to sample the posterior;
it does not preserve the location or rank of an individual forecast member.
The conditional CDF $G_y$ carries the predictive information. Even in
the Gaussian case, this coupling need not agree member by member with
the usual EnKF correction, although both have the same posterior law.

\section{What happens for a finite ensemble?}

Suppose the predictive approximation is
\[
\mu^N=\frac1N\sum_{i=1}^N\delta_{h_i}.
\]
Its exact update is
\[
\pi_y^N=\sum_{i=1}^N w_i\delta_{h_i},\qquad
w_i=\frac{g(y-h_i)}{\sum_{j=1}^N g(y-h_j)}.
\]
The conditional noise law and CDF are
\[
\nu_y^N=\sum_{i=1}^N w_i\delta_{y-h_i},\qquad
G_y^N(u)=\sum_{i=1}^N w_i\ind_{\{y-h_i\le u\}}.
\]
Generating independent noises $\eta_i^f$ and applying the map gives
\[
h_i^a=y-(G_y^N)^{-1}(F_\eta(\eta_i^f)).
\]
Conditional on the fixed forecast ensemble, each output has distribution
$\pi_y^N$. The outputs are independent if fresh independent noises are used.

To see the computation, order the noise locations as
$a_1\le\cdots\le a_N$, where $a_j=y-h_{\sigma(j)}$, and form
$W_j=\sum_{k=1}^j w_{\sigma(k)}$. For $V\sim U(0,1)$, choose
the smallest $j$ with $V\le W_j$. Then $T_y(\eta^f)=a_j$ and the
updated state is $h_{\sigma(j)}$. Tied locations can be combined.
This is inverse-CDF resampling from the weighted posterior.

There are two separate approximations here. First, $\pi_y^N$ updates
the empirical predictive law, which may differ from the true predictive
law. Second, the equal-weight output measure
$N^{-1}\sum_i\delta_{h_i^a}$ has resampling error relative to $\pi_y^N$.
Exact distribution for each draw does not mean exact equality of
these two finite measures. In general, arbitrary posterior weights
cannot be represented by the same number of equal-weight atoms exactly.

\section{A continuous predictive approximation for the SV model}

The Gaussian state transition provides an alternative to replacing the
predictive distribution by Dirac masses. If the previous filtering
approximation is $N^{-1}\sum_i\delta_{h_{t-1}^{a,(i)}}$, then its exact
one-step propagation is the mixture
\[
p_f^N(h)=\frac1N\sum_{i=1}^N
\varphi_q(h-m_i),\qquad
m_i=\mu_h+\phi(h_{t-1}^{a,(i)}-\mu_h),
\]
where $\varphi_q$ is the $N(0,q^2)$ density.

Because $q>0$, this is a continuous positive density. We can use it in
the integral formula for $G_y$ and compute its generalized inverse
numerically. This gives continuous updated particles rather than
resampling only the finitely many forecast locations. It is exact for
the mixture predictive model if the integrals and inverses are exact.
The previous filtering approximation and numerical errors still remain.

\begin{algorithm}[htbp]
\caption{One-step noise-transport update}
\begin{algorithmic}[1]
\Require Predictive law $\mu$; observed centered return $r-\mu_r\ne0$;
ensemble size $N$.
\State $y\gets\log(r-\mu_r)^2$.
\State Construct $G_y$ from \eqref{eq:conditional-cdf}, or its
density form if available.
\For{$i=1,\ldots,N$}
    \State Draw $\eta_i^f\sim g$ independently.
    \State $v_i\gets F_\eta(\eta_i^f)$.
    \State $\eta_i^a\gets G_y^{-1}(v_i)$.
    \State $h_i^a\gets y-\eta_i^a$.
\EndFor
\State \Return $N^{-1}\sum_{i=1}^N\delta_{h_i^a}$.
\end{algorithmic}
\end{algorithm}

At the population level, the theorem specifies an exact posterior
representation. Its practical use still depends on computing $G_y$.
Since $F_\eta(\eta_i^f)$ is uniform, one may equivalently generate
$v_i$ directly from $U(0,1)$. The paired forecast representation is useful
for expressing the update, but it is not required for inverse-CDF sampling.
The construction does not remove the Bayesian conditioning step; it
expresses that step as a nonlinear transformation of known noise.
For this one-dimensional SV model, quadrature and inverse-CDF evaluation
are possible directions to implement and check. No numerical performance
or finite-ensemble convergence rate is claimed in this draft.

\end{document}
